\documentclass[10pt,a4paper]{amsart}
\usepackage[margin=19mm]{geometry}
\usepackage{amsmath,amssymb,mathtools}
\usepackage[scr=rsfs,cal=euler]{mathalfa}
\usepackage{xfrac}
\usepackage[unicode,hidelinks,bookmarks=false]{hyperref}
\newcommand{\ie}{i.e.\ }
\newcommand{\cf}{cf.\ }
\newcommand{\resp}{respectively\ }
\newcommand{\Subsection}{\subsection}
\providecommand{\nobreakdash}{\nobreak\hbox{-}}
\setlength{\emergencystretch}{2em}
\multlinegap0pt
\usepackage{bm}
\usepackage{bbm}
\usepackage{dsfont}
\usepackage{xspace}
\usepackage{cancel}
\usepackage{mathtools}
\usepackage{amsthm}
\makeatletter
\@ifundefined{Theorem}{\newtheorem{Theorem}{Theorem}}{}
\@ifundefined{Lemma}{\newtheorem{Lemma}[Theorem]{Lemma}}{}
\makeatother
\DeclareMathOperator{\dcl}{dcl}
\DeclareMathOperator{\tp}{tp}
\title[Vaught's conjecture for theories of discretely ordered struc]{Vaught's conjecture for theories of discretely ordered structures}
\author{Predrag Tanovi\'c}
\date{Source: arXiv:2212.13605v3 (CC BY 4.0)}
\begin{document}
\maketitle

\noindent\emph{Source and licence.} Vaught's conjecture for theories of discretely ordered structures. Version arXiv:2212.13605v3 (CC BY 4.0), distributed under the Creative Commons Attribution 4.0 International licence, as recorded in the arXiv licence metadata for this exact version and shown on the abstract page https://arxiv.org/abs/2212.13605v3. Full rights notice: licenses/arxiv-2212.13605-CC-BY-4.0.txt.\par\smallskip

\noindent\textit{Editorial scope.} Counted unit: the Lemma whose source label is \texttt{Lemma\_basic\_simple} in the article (source0.tex lines 138--148, source label \texttt{Lemma\_basic\_simple}), delivered together with the article's own complete original proof of it (source0.tex lines 149--172, one \texttt{proof} environment of 24 lines). No mathematics is written, repaired or re-derived: statement and proof are the article's own text line for line. Required context retained uncounted: none. Every definition, display and label of those lines is retained unchanged. Omitted from this package and not redistributed: 1--137, 173--387. Nothing inside a retained excerpt is omitted or altered: every retained body is delivered exactly as the article's lines are written, so no omission receipt of any kind accompanies this package. Citations: the retained excerpts carry no citation command at all, so no bibliography entry of the article is transcribed and no reference is redistributed. Reference adaptation: none was needed and none was made; every reference inside the delivered unit resolves inside it, so no unresolved reference or citation is produced. Printed slips kept literal: no printed word, symbol, hypothesis or number of the counted statement or of its proof was altered. Layout and class: the article's own class is \texttt{amsart} with packages including amsmath, amssymb, amsthm, enumerate, mathabx, geometry; the wrapper is the lane's pinned \texttt{amsart} editorial wrapper at 10pt on A4, which restates the theorem-like environments the retained text needs and adds the article's own macro definitions that the retained text needs (\texttt{\textbackslash dcl}, \texttt{\textbackslash tp}), with extra wrapper packages (bm, bbm, dsfont, xspace, cancel, mathtools). The delivered numbering is the unit's own. Assets: no retained span contains a figure environment, a graphics-inclusion command, a PDF-page inclusion, a table or a TikZ picture; no image file is copied into this package, the package declares no asset dependency and no image file is redistributed with it. Licence: version arXiv:2212.13605v3 of this article is distributed under the Creative Commons Attribution 4.0 International licence, as recorded in the arXiv licence metadata for this exact version and shown on the abstract page https://arxiv.org/abs/2212.13605v3. A full rights notice is delivered at licenses/arxiv-2212.13605-CC-BY-4.0.txt. Characters: every retained excerpt is pure ASCII, so the delivered engine, XeLaTeX with its default Latin Modern text font, reports no missing character.
\begin{Lemma}\label{Lemma_basic_simple}
If   $p\in S_1(A)$ is a simple type, as witnessed by $C\subset\dcl(A)$ and $(D,<)$, then:

(a) $P(M):=C\cup p(M)$ is an initial part  of $(D,<)$ that is type-definable over $A$.

(b) There exists an  $A$-definable, discretely ordered, initial part $D_0\subseteq D$  with $P(M)\subseteq D_0$. In particular, $(p(M),<)$ is a discrete order.  

(c) Every nonempty, relatively $M$-definable subset of  $P(M)$ has a  minimum.  

(d) If $C$ and $(D_1,<_1)$ also witness the simplicity of $p$ such that $<_1$ and $<$ agree on $C$, then they agree on $P(M)$, too. 
\end{Lemma}

\begin{proof}To simplify the notation, assume $A=\emptyset$ and  write $x<y$ instead of $x\in D\land y\in D\land x<y$. 

(a) Since $C$ is an initial part of $(D,<)$,    $C<p(M)$ holds. 
To prove that $P(M)$ is an initial part of $D$, it suffices to show that $a\in p(M)$ and $b< a$ together imply $b\in P(M)$. Let $\phi(x)\in \tp(b)$. 
Then  $\exists y(\phi(y)\land y< x)\in\tp(a)=p$. Since $p$ is   a $C$-type, there are $c\in C$ and $b'\in D$ with $\models\phi(b')\land b'< c$.
Then $b'< c$ implies $b'\in C$, so   $\phi(x)$ is realized by $b'\in C$. We have just shown that every formula $\phi(x)\in \tp(b)$ is realized in $C$, so the type $\tp(b)$ is finitely satisfied in $C$; 
If $\tp(b)$ is algebraic  then $b\in C$; otherwise, $\tp (b)$ is a $C$-type, so $\tp(b)=p$ holds by the uniqueness of $p$.
In both cases we have $b\in P(M)$, so $P(M)$ is an initial part of $D$. Then it is is easy to see that $P(M)$ is defined by  $P(x):=\{\exists y(\theta(y) \land   x<y)\mid \theta(y)\in p(y)\}$. 

\smallskip
(b) Let $c_0=\min C$ and let $\phi(x)$ be a formula expressing that the interval $[c_0,x]$ is discretely ordered by $<$.  For all $c\in C$ the interval $[c_0,c]$ is finite, so $\models  \phi(c)$ holds. Since $p$ is a $C$-type, $\phi(x)\in p$. Let $D_0$ be the union of all intervals $[c_0,a]\subset D$ where $a\in\phi(M)$.  

\smallskip
(c) Let $\emptyset\neq D_1\subset P(M)$  where $D_1$ is relatively defined by $\theta(x,\bar b)$ ($\bar b\in M$). Pick $a\in p(M)$ with $[c_0,a]\cap D_1\neq\emptyset$.  Let $\phi(x,\bar y)$ be a formula expressing that the set
$[c_0,x]\cap\theta(M,\bar y)$ is nonempty and has minimum. 
In particular, $\models\phi(a,\bar b)$ holds. For all $c\in C$ the interval $[c_0,c]$ is finite, so $\models \forall \bar y \left(\exists z(c_0\leq z\leq c\land \theta(z,\bar y))\rightarrow \phi(c,\bar y)\right)$. Since $p$ is a $C$-type and $a\in p(M)$, 
\begin{center} 
$\models \forall \bar y \left(\exists z(c_0\leq z\leq a\land \theta(z,\bar y))\rightarrow \phi(a,\bar y)\right)$.  
\end{center}
In particular,   $\models  \exists z(c_0\leq z\leq a\land \theta(z,\bar b))\rightarrow \phi(a,\bar b)$.  Now  $[c_0,a]\cap D_1\neq \emptyset$ implies $\models  \exists z(c_0\leq z\leq a\land \theta(z,\bar b))$, so $\models\phi(a,\bar b)$ and $\min ([c_0,a]\cap D_1)$ exists. Hence $\min D_1$ exists, too.

\smallskip
(d) Suppose not. Then for all $a\in p(M)$ the set $D(a)$ defined by $x<a\land x>_1 a$ is nonempty so, by part (c), $b=\min_{<} D(a)$ is well-defined. By part (a), $C<_1p(M)$ implies $b\notin C$, so  $b<a$ implies $b\in p(M)$. In particular, $D(b)\neq\emptyset$ and for any $b'\in D(b)$ we have $b'<b<a\land b'>_1b>_1 a$, contradicting the minimality of $b$ in $D(a)$. 
\end{proof}

\end{document}
